\documentclass[../../main.tex]{subfiles}
\begin{document}
\begin{oframed}
  \subfile{problem}
\end{oframed}

{\bf [解答]}

\paragraph{(1)}
$h(x)=6\sin(x/6)$とおく．
\begin{align}
  h'(x)=\cos(x/6)
\end{align}
だから，$P$ ，$Q$における接線$l_p$，$l_q$は，
\begin{align}
  \begin{cases}
    l_p&:y=\dfrac{1}{2}(x-2\pi)+3\sqrt{3}\equiv f(x) \\
    l_q&:y=-(x-6\pi)\equiv g(x)
  \end{cases}
\end{align}
だから，$R$はこれらの交点で，
\begin{align}
  R(14\pi/3-2\sqrt{3},4\pi/3+2\sqrt{3})
\end{align}
である．$\cdots$(答)

\paragraph{(2)}
グラフの概形は\cref{fig:utokyo_1962_6}のようになる．ただし，
\begin{align}
  t=14\pi/3-2\sqrt{3}
\end{align}
である．
\begin{figure}[htb]
  \centering
  \begin{tikzpicture}[x=0.42cm, y=0.55cm, >=stealth,
    every node/.style={font=\small}]
    \def\piVal{3.14159265}
    \def\twoPi{6.283185}
    \def\sixPi{18.849556}
    \def\tVal{11.19688}
    \def\yR{7.65267}
    \def\yP{5.19615}

    \coordinate (O) at (0, 0);
    \coordinate (P) at (\twoPi, \yP);
    \coordinate (R) at (\tVal, \yR);
    \coordinate (Q) at (\sixPi, 0);

    % 斜線領域
    \fill[pattern=north east lines, pattern color=blue!40]
    (P) -- (R) -- (Q)
    -- plot[domain=\sixPi:\twoPi, samples=60] (\x, {6*sin(\x*180/(\piVal*6))})
    -- cycle;

    % 座標軸
    \draw[->, thick, gray!80!black] (-0.8, 0) -- (20.5, 0) node[right] {$x$};
    \draw[->, thick, gray!80!black] (0, -0.8) -- (0, 9.0) node[above] {$y$};
    \node[below left, font=\footnotesize] at (O) {$\mathrm{O}$};

    % 曲線 y = 6 sin(x/6)
    \draw[line width=1.2pt, blue!85!black, domain=0:20.0, samples=80]
    plot (\x, {6*sin(\x*180/(\piVal*6))});

    % 接線 lp, lq
    \draw[line width=1.0pt, red!85!black]
    (\twoPi - 2.5, {\yP - 2.5*0.5}) -- (\tVal + 0.8, {\yR + 0.8*0.5})
    node[above left, font=\footnotesize] {$l_p$};
    \draw[line width=1.0pt, teal!85!black]
    (\tVal - 0.8, {\yR + 0.8}) -- (\sixPi + 1.2, {-1.2})
    node[below right, font=\footnotesize] {$l_q$};

    % 垂線
    \draw[densely dotted, thick, gray!70!black] (\twoPi, 0) node[below, font=\footnotesize] {$2\pi$}
    -- (P)
    -- (0, \yP) node[left, font=\footnotesize] {$3\sqrt{3}$};
    \draw[densely dotted, thick, gray!70!black] (\tVal, 0) node[below, font=\footnotesize] {$t$}
    -- (R);
    \node[below, font=\footnotesize] at (\sixPi, 0) {$6\pi$};

    % 代表点
    \node[circle, fill=blue!85!black, inner sep=1.3pt] at (P) {};
    \node[above left=1pt, font=\footnotesize] at (P) {$P$};
    \node[circle, fill=red!85!black, inner sep=1.3pt] at (R) {};
    \node[above right=1pt, font=\footnotesize] at (R) {$R$};
    \node[circle, fill=blue!85!black, inner sep=1.3pt] at (Q) {};
    \node[above right=2pt, font=\footnotesize] at (Q) {$Q$};

    % 曲線ラベル
    \node[above, font=\footnotesize, blue!85!black] at (3.5*\piVal, 6.0) {$y=6\sin\frac{x}{6}$};
  \end{tikzpicture}
  \caption{$y=6\sin(x/6)$ と接線 $l_p, l_q$ および囲まれた領域}
  \label{fig:utokyo_1962_6}
\end{figure}
従って，求める面積$S$として，
\begin{align}
  S&=\int_{2\pi}^tfdx+\int_t^{6\pi}gdx-\int_{2\pi}^{6\pi}hdx \nonumber\\
  &=\dfrac{8}{3}\pi^2+8\sqrt{3}\pi-63\nonumber
\end{align}
である．$\cdots$(答)
\newpage
\end{document}
